\subsubsection{A division-free theta heat-loss formula}
\label{subsec:thetaspectralloss}

A spatially dependent comparison
$k_nU_{n+1}\le c_nU_n$ is difficult to evaluate
where $U_n$ becomes small. The theta convolution
admits another exact representation in which
neither the coefficient $c_n$ nor division by
$U_n$ appears. It also exposes the obstacle
to proving noncentral positivity by replacing
cosine transforms with positive moments.

Put $T=1/2$, $t=T-s$, $0\le s\le T$, and
$k_n=(n+1)(2n+1)$. Write
$h_T(u)=\Phi(|u|)e^{Tu^2}/2$, so that
$H_T(z)=\int_{\mathbb R}h_T(u)e^{izu}\,du$.
For $n\ge1$, define the strictly nonnegative
pair weight
\begin{equation}
W_n(u,v)=\frac{h_T(u)h_T(v)}{(2n)!}
                  (u-v)^{2n}.
\label{eq:heatpairweight}
\end{equation}
Here $P_s=e^{s\partial_x^2/2}$ denotes the
forward Gaussian heat semigroup.

\begin{theorem}[Exact theta heat-loss identity]
\label{thm:thetaspectralloss}
For $s\in[0,T]$, $n\ge1$, and real $x$,
\begin{align}
U_n(s,x)
&=\iint_{\mathbb R^2}W_n(u,v)
e^{-s(u+v)^2/2}e^{-s(u-v)^2/2}\nonumber\\
&\hspace{25mm}\cdot\cos(x(u+v))\,du\,dv,\label{eq:heatpairU}\\
P_sU_n(0,x)
&=\iint_{\mathbb R^2}W_n(u,v)
e^{-s(u+v)^2/2}\cos(x(u+v))\,du\,dv.
\label{eq:heatpairP}
\end{align}
Consequently, with $\Delta_n=P_sU_n(0)-U_n(s)$,
\begin{align}
\Delta_n(s,x)=
\iint_{\mathbb R^2}W_n(u,v)
e^{-s(u+v)^2/2}
\bigl(1-e^{-s(u-v)^2/2}\bigr)\nonumber\\
\hspace{28mm}\cdot\cos(x(u+v))\,du\,dv.
\label{eq:thetaloss}
\end{align}
For $s>0$, $\Delta_n(s,0)>0$. In addition,
put
\begin{equation}
M_{n+1}(s)=
\iint_{\mathbb R^2}
\frac{h_T(u)h_T(v)(u-v)^{2n+2}}{(2n+2)!}
e^{-s(u+v)^2/2}\,du\,dv>0.
\end{equation}
Then the following bound is valid without
assuming any sign for $U_n(s,x)$:
\begin{equation}
\boxed{|\Delta_n(s,x)|\le
s\,k_nM_{n+1}(s),\qquad x\in\mathbb R.}
\label{eq:thetaheatlossbudget}
\end{equation}
\end{theorem}

\begin{proof}
The coefficient of $y^{2n}$ in
$H_t(x+iy)H_t(x-iy)$ equals
$(2n)!^{-1}$ times the Fourier transform
of $h_t(u)h_t(v)(u-v)^{2n}$ at $u+v$.
Since $u^2+v^2=((u+v)^2+(u-v)^2)/2$,
replacing $h_T$ with $h_{T-s}$ gives
\eqref{eq:heatpairU}. Gaussian convolution
in the $x$ variable multiplies the
$(u+v)$ Fourier frequency by
$e^{-s(u+v)^2/2}$, proving
\eqref{eq:heatpairP}. Subtraction proves
\eqref{eq:thetaloss}.

At $x=0$ the integrand of
\eqref{eq:thetaloss} is nonnegative,
and is strictly positive away from the
diagonal $u=v$ for $s>0$.
The theta kernel is positive on a
set of full measure, giving strictness.
Finally, $0\le1-e^{-a}\le a$ for $a\ge0$,
$|\cos|\le1$, and
$(2n+2)!/(2(2n)!)=k_n$ prove
\eqref{eq:thetaheatlossbudget}.
\end{proof}

\begin{corollary}[An all-order spatially decaying
division-free estimate]\label{cor:spectrallossdecay}
For fixed $s>0$, let $K_{n,s}(w)$ be the
push-forward of the nonnegative measure
in \eqref{eq:thetaloss} under $w=u+v$.
Then $K_{n,s}$ is even, nonnegative
and rapidly decreasing, and
\[
\Delta_n(s,x)=\int_{\mathbb R}K_{n,s}(w)
                      \cos(xw)\,dw.
\]
For every integer $j\ge1$ and $x\ne0$,
\begin{equation}
|\Delta_n(s,x)|\le
\min\left\{
s k_nM_{n+1}(s),
\frac{\|K_{n,s}^{(2j)}\|_{L^1}}{|x|^{2j}}
\right\}.
\label{eq:lossderivatives}
\end{equation}
In particular the explicit inequality
\begin{equation}
P_sU_n(0,x)>
\min\left\{s k_nM_{n+1}(s),
\frac{\|K_{n,s}''\|_1}{x^2}\right\}
\label{eq:noncircularband}
\end{equation}
implies $U_n(s,x)>0$, with no assumed
positivity of the unknown $U_n(s,x)$.
\end{corollary}
\begin{proof}
The Jacobi transformation makes $h_T$
even and smooth; its superexponential
tails give all required Schwartz bounds.
Integrating out $u-v$ preserves smoothness
and rapid decrease in $w$.
Repeated integration by parts in the
cosine transform gives the derivative
estimates. The final implication follows
from $U_n=P_sU_n(0)-\Delta_n$.
\end{proof}

\begin{proposition}[An exact noncentral sign obstruction]
\label{prop:signedlossatom}
A positive even spectral weight does not force
$\Delta_n(s,x)\ge0$ away from $x=0$,
even inside $0<s<1/2$.
Let the positive symmetric atomic measure
have masses $\mu(0)=1$,
$\mu(-1)=\mu(1)=1/10$; use the same
pair formula with $h_T(u)\,du$ replaced
by $d\mu(u)$. At the admissible time
$s=2\log(5/4)<1/2$ one has
\[
\boxed{\Delta_1(s,0)=\frac{869}{15625}>0,
\qquad \Delta_1(s,\pi)=-\frac{131}{15625}<0.}
\]
\end{proposition}
\begin{proof}
Write $r=e^{-s/2}=4/5$. The terms
with one atom at $0$ and the other at
$\pm1$ contribute, before division
by $2!=2$, $\pm4(1/10)r(1-r)$,
with the negative sign at $x=\pi$.
The two opposite nonzero atoms
contribute
$2(1/10)^2(2)^2(1-r^4)=738/15625$.
Hence at $x=\pi$ the sum is
$\frac12(-8/125+738/15625)
=-131/15625$, whereas at $x=0$
it equals $869/15625$.
The time inequality follows from
$\log(1+z)<z$ for $z>0$.
This is a counterexample to an
unjustified sign inference, not a
claim about a zero of the Riemann
transform.
\end{proof}

\paragraph{A weighted Duhamel alternative.}
For any nonnegative locally bounded
spatial coefficient $c_n(s,x)$ put
$D_n=k_nU_{n+1}-c_nU_n$ and let
$\mathcal Q_{c_n}(s,r)$ denote the
positive evolution operator of
$\frac12\partial_x^2-c_n$.
Duhamel's identity is
\begin{equation}
U_n(s)=\mathcal Q_{c_n}(s,0)U_n(0)
-\int_0^s\mathcal Q_{c_n}(s,r)D_n(r)\,dr.
\label{eq:weighteddivisionfree}
\end{equation}
No division by $U_n$ is needed.
The pointwise condition $D_n\le0$
suffices, but is stronger than the
integrated condition that the positive
part of the Duhamel source be dominated
by the initial semigroup term. Such
an integrated estimate could tolerate
localized violations of the proposed
spatial comparison. It remains
unproved for the genuine theta kernel.

The nonoscillatory budget
\eqref{eq:thetaheatlossbudget} is positive
and independent of Fourier height,
whereas $P_sU_n(0,x)$ decays to zero
as $|x|\to\infty$. Therefore that
budget alone cannot prove
\eqref{eq:noncircularband} for
arbitrarily large $|x|$. The derivative
bounds \eqref{eq:lossderivatives}
supply spatial decay but no
theta-specific \emph{relative} lower
bound on the positive baseline.
This is the precise remaining
Fourier-sign problem for this method.
The algebraic normalizations and
the finite atomic sign obstruction are
checked with exact rational arithmetic
in \path{certificates/xi_division_free_spectral_loss_exact.py}
and at the Lean algebra level in
\path{XiKernel/ThetaSpectralLoss.lean}.
Neither check formalizes the complete
theta Fourier integrals or proves RH.
